El crecimiento acumulado de los principales indicadores monetarios y financieros durante el período 2015–2025 evidencia que las remesas familiares fueron la variable con mayor expansión, registrando un incremento aproximado del
371.7%
respecto al período base. En segundo lugar se ubicó el numerario en circulación con un crecimiento cercano al
313.7%
, mientras que los medios de pago (M2), el crédito bancario al sector privado y el Producto Interno Bruto (PIB) mostraron incrementos más moderados, comprendidos entre el
111.1%
y el
181.5%
.
Estos resultados reflejan que, durante la última década, las remesas familiares se consolidaron como uno de los principales motores del crecimiento económico y de la expansión de la liquidez en Guatemala. Asimismo, el importante aumento del numerario en circulación sugiere que el efectivo continúa desempeñando un papel predominante dentro de las transacciones económicas del país, situación que puede estar asociada con una elevada utilización del dinero en efectivo y con los desafíos que aún existen en materia de inclusión financiera.
En conjunto, el comportamiento de estos indicadores muestra que el crecimiento del sistema monetario no ha sido homogéneo entre todas las variables. Mientras el PIB y el crédito bancario mantuvieron una evolución más estable y consistente con el crecimiento económico de largo plazo, las remesas familiares y el numerario fueron los componentes que registraron la mayor dinámica de expansión, resaltando su importancia dentro de la economía guatemalteca y su potencial influencia sobre la liquidez y el consumo interno.